% !TeX root = main.tex
\chapter*{АНОТАЦІЯ}
\phantomsection
\addcontentsline{toc}{chapter}{АНОТАЦІЯ}

Денисюк Д.~О. Інформаційна технологія виявлення ознак аномального вмісту в мультимедійних даних вебсередовища. Кваліфікаційна наукова праця на правах рукопису.

Дисертація на здобуття наукового ступеня доктора філософії за спеціальністю 126 «Інформаційні системи та технології» галузі знань 12 «Інформаційні технології». Хмельницький національний університет, Хмельницький, 2026.

Дисертацію присвячено розв’язанню науково-прикладного завдання щодо формування обґрунтованого рішення про ознаки аномального вмісту в мультимедійних даних за неповних і різнорідних спостережень під час їх оброблення у вебсистемі. Формальна відповідність заявленому формату й зорова правдоподібність не виключають порушень структури, нетипових спектральних і зорових ознак або відхилень у перебігу подій і поведінці компонентів. Аналіз окремого представлення не дає змоги простежити зв’язок між властивостями носія, операцією оброблення та можливим наслідком. У роботі термін «аномальний вміст» позначає дані або структурні значення, що відхиляються від установлених умов формату, статистичного опису чи контрольованого маршруту оброблення. Спостережуване відхилення не ототожнюється з доведеною загрозою, виконуваним або шкідливим кодом.

Метою дослідження є формування інформаційної технології виявлення визначених структурних, спектральних і зорових ознак, а також ознак подій і поведінки, що характеризують аномальний вміст у мультимедійних даних вебсередовища, на основі моделі мультимодального представлення прихованої загрози, методу аналізу спектральних, зорових і структурних ознак та методу поведінкового узгодження ознак різного походження. Об’єктом дослідження є процес формування рішення щодо ознак аномального вмісту під час приймання й оброблення мультимедійних даних компонентами вебсистеми. Предмет дослідження становлять модель мультимодального представлення прихованої загрози, метод аналізу спектральних, зорових і структурних ознак, метод поведінкового узгодження ознак різного походження та інформаційна технологія виявлення ознак аномального вмісту в мультимедійних даних вебсередовища.

Для досягнення мети використано системний аналіз, математичне моделювання, цифрове оброблення зображень і відео, структурний аналіз мультимедійних контейнерів, методи машинного навчання, описову статистику, групове повторне вибирання та парне зіставлення конфігурацій. Системний аналіз і математичне моделювання застосовано до формалізації маршруту об’єкта, мультимодального представлення та правила прийняття рішення, а цифрове оброблення і структурний аналіз --- до формування статичних ознак. Статистичні методи використано для оцінювання матриць помилок, невизначеності показників і відмінностей між конфігураціями.

Наукова новизна полягає в розробленні моделі мультимодального представлення прихованої загрози, у якій мультимедійний об’єкт подано як багаторівневий інформаційний контейнер. Незмінне байтове подання, контрольовано декодований вміст, структурна карта, супровідні відомості і причинно пов’язаний журнал подій зберігаються як окремі сутності зі спільним ідентифікатором і відомостями про походження. Структурні, спектральні й зорові ознаки, а також ознаки подій і поведінки формально пов’язуються у представленні $Z$ зі збереженням часткових оцінок, локалізацій і маски доступності. Така побудова дає змогу розмежувати оцінку ризику, упевненість, аналітичний стан і технологічну дію та не трактувати недоступну групу ознак як підтвердження нормального стану. Повне представлення $Z$ із причинно пов’язаними подіями експериментально не формувалося.

Удосконалено метод аналізу спектральних, зорових і структурних ознак мультимедійних даних. Гілка спектрального та зорового аналізу формує високочастотні залишки, характеристики дискретного косинусного й дискретного хвилькового перетворень, ентропійні ознаки та часові агрегати кадрів. Структурна гілка перевіряє заголовки, межі елементів контейнера, службові поля, завершальні маркери й додаткові байтові області. Окремі оцінки та локалізації двох гілок зберігаються до їх інтеграції. Скорочені числові проєкції первинно оцінено для PNG і MP4; перенесення незмінних параметрів на JPEG, WebP і WebM перевірено без повторного навчання та добору порогів. Повний склад статичного опису й локалізації експериментально не оцінено.

Розроблено метод поведінкового узгодження ознак різного походження, який передбачає зіставлення статичного опису з подіями оброблення того самого об’єкта за ідентифікатором, часовими межами, фазою та причинним зв’зком. Метод визначає вхідний і вихідний контракти, причинно-фазовий допуск, кодування подій, кероване об’єднання, граничні стани та умови перевірки. Для експериментального оцінювання реалізовано скорочену модель узгодження статичних оцінок з урахуванням формату. Вона опрацьовує вже сформовані оцінки першого методу й не є реалізацією повного модуля кодування подій, формування ознак подій і поведінки та перехресного зіставлення.

Набула подальшого розвитку інформаційна технологія, у якій модель і два методи організовано в п’ять етапів: попереднє оброблення, формування ознак, інтеграція ознак, прийняття рішення та реєстрація результатів. Початковий об’єкт зберігається незмінним, реконструйований об’єкт проходить окремий цикл, а типізований вихід містить версії складників, локалізації, маску доступності та посилання на журнал. Аналітичний результат відокремлено від визначеної політикою дії: допуску, експертної перевірки, реконструкції або блокування. Локально перевірено переходи $S_1$--$S_5$, атомарну реєстрацію, повтор, відновлення й безпечне завершення за двох технічних відмов. Практичне значення одержаних результатів полягає у визначенні структури простежуваного технологічного контуру між каналом приймання мультимедійного об’єкта та прикладною логікою вебсистеми, а також складу машинозчитуваного запису для повторної перевірки.

Експериментальне дослідження охопило шість серій: ретроспективну, внутрішню групово відокремлену та чотири зовнішні, доповнені окремим ресурсним випробуванням і локальним аудитом технологічних станів. За збереженим підсумковим описом внутрішня серія містила $200$ первинних груп PNG і MP4 та $400$ похідних об’єктів із поділом $60:20:20$. На перевірній частині з $80$ об’єктів спільна числова проєкція $Z_c$ дала $40$ правильно визначених контрольних об’єктів, жодного хибнопозитивного рішення, два хибнонегативні рішення та $38$ правильно визначених змінених об’єктів. Із матриці повторно обчислюються частка правильних рішень $0{,}9750$, повнота $0{,}9500$ і $F_1=0{,}9744$; повідомлену площу під характеристичною кривою $0{,}9869$ без неперервних пооб’єктних оцінок незалежно не перевірено.

Чотири зовнішні набори містили $38$ реальних відеозаписів із $12$ пристроїв і $24$ груп походження. Перша зовнішня серія на $56$ похідних MP4 не підтвердила переваги спільної проєкції над структурною конфігурацією: збалансована частка правильних рішень становила відповідно $0{,}7857$ і $0{,}8333$. У багатоформатній серії $500$ базових фрагментів утворили $10\,000$ похідних об’єктів форматів PNG, JPEG, WebP, MP4 і WebM. Для структурної конфігурації, конфігурації спектральних і зорових ознак та спільної конфігурації одержано збалансовані частки правильних рішень $0{,}7636$, $0{,}5001$ і $0{,}7916$. Кількість похідних об’єктів не прирівнювалася до кількості незалежних джерел.

Скорочену модель узгодження статичних оцінок перевірено окремо на $2\,000$ об’єктах, утворених зі $100$ базових фрагментів восьми нових записів чотирьох пристроїв. Після фіксації моделі й порога збалансована частка правильних рішень змінилася від $0{,}7827$ для спільної конфігурації до $0{,}8307$ для моделі узгодження. Модель усунула $76$ хибнопозитивних рішень, але додала $84$ хибнонегативні рішення та не виявила $500$ об’єктів лише зі спектральною зміною. Точний двобічний критерій Мак-Немара для звичайної частки правильних рішень дав $p=0{,}5801$. Отже, встановлено зміну співвідношення між правильним визначенням контрольних об’єктів і повнотою, а не безумовну перевагу моделі.

Ретроспективне оцінювання $110$ об’єктів підтвердило арифметичну узгодженість агрегованої матриці рішень і похідних показників; площу під характеристичною кривою не визначено через відсутність неперервних пооб’єктних оцінок. Середній час чотирьох зафіксованих операцій становив $418{,}09$~мс, але не характеризує повний п’ятиетапний цикл. У п’яти ресурсних запусках 95-й процентиль часу для структурної, спільної та узгоджувальної конфігурацій становив відповідно $2{,}136$, $853{,}971$ і $1078{,}904$~мс. У локальному аудиті $35$ циклів завершилися цілісним $S_5$, а десять очікуваних відмов --- без $S_5$ та без зовнішньої дії; один чотирипроцесний запуск дав $22{,}827$ об’єкта за секунду. Рейтинг $K_I$ використано лише як вторинний засіб зіставлення достовірності, часових витрат, пам’яті та стійкості між форматами.

Межі висновків зумовлені залежністю похідних об’єктів від спільних базових фрагментів і обмеженою кількістю незалежних пристроїв. Повторне виділення ознак із тих самих $56$ і $2\,000$ файлів MP4 зберегло збалансовану правильність структурної гілки $0{,}8333$, але змінило її для спільної проєкції від $0{,}7857$ до $0{,}6905$ і від $0{,}7987$ до $0{,}6520$. Тому поточну еталонну реалізацію модуля виділення спектральних ознак не ототожнено з відсутньою історичною реалізацією.

На $1\,500$ допустимо змінених контрольних об’єктах частка хибнопозитивних структурних рішень для приватних розширень або відступів становила $1{,}0000$ у кожному форматі; отже, структурний результат не є самодостатнім доказом прихованого впливу. У парному розвідувальному аналізі кожне з $2\,500$ форматних представлень із внесеною спектральною зміною зберегло ненульову відмінність після декодування, однак через відсутність просторово-часових масок, збережених до кодування, первинну локалізацію зміни не оцінено.

Причинно пов’язаних журналів і повного представлення ознак подій і поведінки не сформовано, а доступна технічна ознака не спрацювала у $4\,832$ фізично окремих похідних відеооб’єктах; цю кількість не прирівнюють до кількості незалежних джерел. Експерименти не охоплювали повний метод поведінкового узгодження, виявлення виконуваного шкідливого коду, реконструкцію, виконання зовнішніх приписів, розподілену систему та тривале промислове навантаження. Одержані результати підтверджують статичну основу першого методу та локальну прототипну реалізацію технологічних переходів у межах перевірених даних, але не обґрунтовують безумовного автоматичного блокування.

За темою дисертації опубліковано дев’ять наукових праць: дві статті у фахових наукових виданнях України містять основні результати, сім праць засвідчують апробацію матеріалів.

\textbf{Ключові слова:} аномальний вміст, мультимедійні дані, структурний аналіз, аналіз спектральних і зорових ознак, мультимодальне представлення, інформаційна технологія, вебсередовище.

\clearpage
